\documentclass[a4paper]{article}
\usepackage[pdftex]{graphicx}
\usepackage[utf8]{inputenc}
\usepackage{enumerate}
\usepackage{icomma}
\usepackage{amssymb}
\usepackage{amsmath}
\usepackage{tikz}
\usepackage{href-ul}
\hypersetup{
	colorlinks=true,
	linkcolor=blue,
	urlcolor=blue}
\usepackage{geometry}
\geometry{a4paper, top=15mm, left=15mm, right=15mm, bottom=15mm,
	headsep=10mm, footskip=12mm}

\begin{document}
	
	\begin{tikzpicture}(10,3)
		\draw[ultra thick](2,0) --(10,0) -- (10,1.3) --(2,1.3) -- (2,0);
		\draw[fill=black](2,0)-- (0,.6) -- (2,.6) -- (2,0);
		\node at (6,.6) {\large Solution sheet from \href{http://www.okuyakl.com}{www.okuyakl.com}};
	\end{tikzpicture}
	\vspace{0.5 cm}
	
	
	\noindent{\bf Task 1.a)}\\
	$$
	\renewcommand{\arraystretch}{2}
	\begin{array}{rcll}
		\tan^2 x &=& {1 \over 4(1-\sin^2 x)} \\
		\tan^2 x &=& {1 \over 4\cos^2 x} \\
		{\sin^2 x \over \cos^2 x} &=& {1 \over 4\cos^2 x} &| \cdot \cos^2x\\
		\sin^2 x &=& { 1 \over 4} &|\sqrt{\quad}\\
		\sin x &=& \pm { 1 \over 2} \\
		x_1 &=& {\pi \over 6}\\
		x_2 &=& -{\pi \over 6} \\
		x_2 &=& -{\pi \over 6} + 2\pi &={11 \pi \over 6} \\
		x_3 &=&\pi-{\pi \over 6} &= {5\pi \over 6}\\
		x_4 &=& \pi -\left(-{\pi \over 6}\right) &= {7 \pi \over 6}\\
		
	\end{array}
	$$
	
	\noindent{\bf Task 1.b)}\\
	$$
	\renewcommand{\arraystretch}{2}
	\begin{array}{rcl}
		\sin x \cos x &=& \sqrt{1-\sin^2 x}\\
		\sin x \cos x &=& \sqrt{\cos^2 x} \\
		\sin x \cos x &=& \cos x \\
		\sin x &=& 1 \quad \lor \quad \cos x =0\\
		x_1 &=& {\pi \over 2} \\
		x_2 &=& {3 \pi \over 2}
	\end{array}
	$$
	
	\noindent{\bf Task 1.c)}\\
	$$
	\renewcommand{\arraystretch}{2}
	\begin{array}{rcll}
		{\cos^2 x \over \sin^2 x} &=& {3 \over 4 \sin^2 x}\\
		\cos^2 x &=& {3 \over 4} &|\sqrt{\quad}\\
		\cos x &=& \pm \sqrt{3 \over 4} \\
		x_1 &=& {\pi \over 6} \\
		x_2 &=& \pi-{\pi \over 6} \\
		x_2 &=& {5\pi \over 6} \\
		2\pi - x_3 &=& {\pi \over 6} \\
		x_3 &=& { 11 \pi \over 6} \\
		2\pi - x_4 &=& {5\pi \over 6} \\
		x_4 &=& { 7 \pi \over 6}
	\end{array}
	$$
	
	\vspace{2cm}
	
	\noindent{\bf task 1.d)}\\
	$$
	\renewcommand{\arraystretch}{2}
	\begin{array}{rcll}
		\sin x + {1 \over 2} &=& { 1 \over 2 \sin x} &|\cdot \sin x\\
		\sin^2 x + {1 \over 2}\sin x &=& { 1 \over 2} & | - {1 \over 2}\\
		\sin^2 x + {1 \over 2}\sin x - { 1 \over 2} &=& 0 &| \sin x = u \\
		u^2 + {1 \over 2} u - {1 \over 2} &=& 0 &\textnormal{midnight formula}\\
		\Rightarrow u_1 &=& -1 \\
		u_2 &=& {1 \over 2} \\
		\sin x &=& -1 \\
		x_1 &=& {3 \pi \over 2}\\
		\sin x &=& {1 \over 2} \\
		x_2 &=& {\pi \over 6} \\
		x_3 &=& {5 \pi \over 6}
	\end{array}
	$$
	
	\noindent{\bf task 1.e)}\\
	$$
	\renewcommand{\arraystretch}{1.5}
	\begin{array}{rcll}
		\sin^2 x + {1 \over 2}\sqrt{2}\sin x + \cos x &=& \cos x &|-\cos x\\
		\sin^2 x + {1 \over 2}\sqrt{2}\sin x &=& 0\\
		\sin x \left(\sin x +{1 \over 2}\sqrt{2}\right) &=&0\\
		\sin x &=& 0 \\
		x_1 &=& 0 \\
		x_2 &=& \pi \\
		x_3 &=&2 \pi \\
		\sin x +{1 \over 2}\sqrt{2}&=&0\\
		\sin x&=& - {1 \over 2}\sqrt{2} \\
		x_4 &=& -{\pi \over 4} + 2 \pi &= {7 \pi \over 4}\\
		x_5 &=& \pi - \left( -{\pi \over 4}\right) &= {5\pi \over 4}
	\end{array}
	$$
	
	\noindent{\bf Task 2.}\\
	a) $$-\sin x+ \sin x =0$$
	b) $$2\cos x - \cos x =\cos x$$
	c) $$\cos x + \cos x = 2 \cos x$$
	d) $$ \sin x - \sin x =0$$
	e)$$\cos x - \cos x=0$$
	f) $$\sin x + \sin x = 2 \sin x$$
	g)$$-\sin x + \sin x =0$$
	
	
	\onecolumn
	\noindent{\bf Task 3.}\\
	a) $$ (\sin^2 x + \cos^2 x)^2 = 1^2 =1$$
	b)$${{ \sin x \over \cos x}\over \sqrt{1+ { \sin^2 x \over \cos^2 x}}}={\sin x \over \sqrt{\cos ^2 x (1+ { \sin^2 x \over \cos^2 x})}}={\sin x \over \sqrt{\cos^2 x + \sin^2 x}}= {\sin x \over 1}=\sin x$$
	
	\noindent{\bf Task 4.}\\
	Further results in the interval [$- \pi ; 5 \pi$] we get by illustrating the sine function:
	
	\includegraphics[width=12 cm]{sin144}
	
	\noindent
	Due to the symmetry of the sine function:
	$$\sin{x}= \sin{(\pi - x)}$$
	So another solution to the above equation is:
	$$x_m= \pi + {\pi \over 5}= 1.2 \pi$$
	For integers $k \in \mathbb{Z}$ the following applies due to the periodicity of the sine function:
	$$ \sin{x}= \sin{(x+2k\pi)}$$
	So now there are all solutions
	$$
	\begin{array}{c|c|c|c|c}
		k & -1 & 0 & 1 & 2 \\
		\hline
		x_n & - & -0.2 \pi & 1.8 \pi & 3.8 \pi \\
		x_m & -0.8\pi & 1.2 \pi & 3.2 \pi & - \\
	\end{array}
	$$
	\noindent {\bf task 5.}\\
	
	\begin{minipage}{0.5\textwidth}
		We first transform the expression in sine so that the shift in the x direction can be seen
		becomes:
		$$ f(x)=2\sin{\left({1 \over 2} \left( x - {\pi \over 2}\right)\right)}$$
	\end{minipage}
	\begin{minipage}{0.5\textwidth}
		\includegraphics[width=7 cm]{f2sin144}
	\end{minipage}
	
	\begin{center}
		\includegraphics[width=7 cm]{../../viecher/eendcomic.pdf}
		
		Here you can go back to the \href{https://www.okuyakl.de/math/m10gozL025/ae025.pdf}{task sheet}
	\end{center}
	
\end{document}